% =====================================================================
%  A Quantum Transformer World Model for Model-Based Reinforcement
%  Learning
%
%  V. Liutarevich
%  Draft v4 --- prepared for Quantum (quantum-journal.org)
%
%  ---------------------------------------------------------------------
%  DOCUMENT CLASS
%  ---------------------------------------------------------------------
%  Quantum requires the `quantumarticle' class. It is in TeX Live but was
%  not installed where this draft was compiled, so the file defaults to
%  `article' and compiles anywhere. To switch for submission, replace the
%  \documentclass line with:
%
%      \documentclass[twocolumn,superscriptaddress]{quantumarticle}
%
%  and delete the \titleblockfallback definitions marked below. Nothing
%  else needs to change: no quantumarticle-specific macros are used in
%  the body.
%
%  ---------------------------------------------------------------------
%  FRAMING (v4 vs v3)
%  ---------------------------------------------------------------------
%  v3 was organised around what the architecture fails to do. That was
%  honest but it was also bad physics writing, because it treated an
%  outcome that expressivity theory PREDICTS as though it were a
%  disappointment. Parity with a capacity-matched classical model at
%  six qubits is not a negative result; it is the result Schuld, Sweke
%  and Meyer's Fourier analysis requires, and observing anything else
%  would have warranted suspicion before celebration. Section 8 now
%  makes that argument explicitly and quantitatively.
%
%  v4 therefore leads with the construction and with what it measures:
%  a working quantum transformer world model, a 26% improvement in
%  open-loop dynamics prediction under partial observation, a decoupling
%  between world-model accuracy and agent return that is a finding about
%  model-based RL in general, hardware execution that transfers, and a
%  precise statement of where the classical-simulability bound stops
%  applying. Nothing is withdrawn or softened --- the limitations
%  section and Appendix C are complete --- but the paper is about an
%  instrument rather than about a scoreboard.
%
%  ---------------------------------------------------------------------
%  PLACEHOLDERS
%  ---------------------------------------------------------------------
%     \pend{...}        inline, for a missing number
%     \pendpara{...}    a paragraph-sized gap
%     pendingbox env    a table or section awaiting the analysis run
%
%  \draftfalse hides all of them. Appendix A lists them in priority
%  order. Everything pending is blocked on one re-evaluation pass over
%  the completed runs; no further training is required.
% =====================================================================
\documentclass[11pt,onecolumn]{article}

\usepackage[margin=1in]{geometry}
\usepackage[T1]{fontenc}
\usepackage[utf8]{inputenc}
\usepackage{amsmath,amssymb}
\usepackage{booktabs}
\usepackage{graphicx}
\usepackage{xcolor}
\usepackage{caption}
\usepackage[colorlinks=true,linkcolor=NavyBlue,citecolor=NavyBlue,
            urlcolor=NavyBlue]{hyperref}
\usepackage[numbers,sort&compress]{natbib}

\definecolor{NavyBlue}{HTML}{1F3864}
\definecolor{PendBG}{HTML}{FFF4E5}
\definecolor{PendFG}{HTML}{8A4B00}

% ---- placeholder machinery -------------------------------------------
\newif\ifdraft
\draftfalse                    % \draftfalse to hide every placeholder

\newcommand{\pend}[1]{%
  \ifdraft\textcolor{PendFG}{\textbf{[}\textit{#1}\textbf{]}}\fi}

\newcommand{\pendpara}[1]{%
  \ifdraft\par\vspace{2pt}\noindent
  \fcolorbox{PendFG}{PendBG}{%
    \begin{minipage}{0.97\linewidth}\small\textcolor{PendFG}{%
      \textbf{PENDING.} #1}\end{minipage}}\par\vspace{4pt}\fi}

\newsavebox{\pendbox}
\newenvironment{pendingbox}[1]
  {\par\vspace{3pt}\noindent
   \begin{lrbox}{\pendbox}\begin{minipage}{0.93\linewidth}\small
   \textcolor{PendFG}{\textbf{AWAITING ANALYSIS --- #1}}\par\vspace{3pt}
   \color{PendFG}}
  {\end{minipage}\end{lrbox}%
   \ifdraft\fcolorbox{PendFG}{PendBG}{\usebox{\pendbox}}\fi
   \par\vspace{4pt}}

% ---- delete the next three lines when switching to quantumarticle ----
\newcommand{\affiliation}[1]{\\[2pt]\small\textit{#1}}
\newcommand{\email}[1]{\\[2pt]\small\texttt{#1}}
\newcommand{\maketitlefallback}{\maketitle}


% ---- figures -----------------------------------------------------------
% Guarded so the document compiles before make_figures.py has been run.
% \detokenize is needed because the filenames contain underscores, which
% are not valid in text mode without it.
\usepackage{etoolbox}
\newcommand{\figorplaceholder}[4]{%
  \IfFileExists{figures/#1}{%
    \begin{figure}[htb]\centering
      \includegraphics[width=#2\linewidth]{figures/#1}
      \caption{#3}\label{fig:#4}
    \end{figure}}{%
    \ifdraft\begin{figure}[htb]\centering
      \fcolorbox{PendFG}{PendBG}{\begin{minipage}{0.88\linewidth}\small
        \textcolor{PendFG}{\textbf{FIGURE NOT BUILT:}
        \texttt{figures/\detokenize{#1}}\\[2pt]
        Run \texttt{python Paper\_B/make\_figures.py}}\end{minipage}}
      \caption{#3}\label{fig:#4}
    \end{figure}\fi}}

\title{\bfseries A Quantum Transformer World Model for\\
       Model-Based Reinforcement Learning}
\author{V. Liutarevich
        \affiliation{\pend{affiliation}}
        \email{\pend{contact email}}}
\date{\today}

\begin{document}
\maketitlefallback

\begin{abstract}
\noindent
Quantum machine learning for reinforcement learning has concentrated on
the policy and the value function. We argue that the better-motivated
target is the \emph{world model}, and we build one. We present a quantum
transformer world model: a transformer over a window of recent latent
states, with a parameterised quantum circuit as the token-mixing block
of one encoder layer, serving as the dynamics prior of a recurrent
state-space model and trained end to end inside a model-based agent. To
our knowledge this is the first quantum model placed in the dynamics
component of a model-based agent. Because latent imagination rolls the
model forward on the prior alone, without observations, the policy is
trained entirely on trajectories generated through a quantum circuit.

The architecture makes a falsifiable prediction, and we test it. A prior
that attends over recent recurrent states should help when the dynamics
are history-dependent and not otherwise. On CartPole with both velocity
components hidden, the quantum transformer prior outperforms its
capacity-matched classical twin by $144$ return points over eight seeds
($353.9$ against $209.9$, a $69\%$ improvement; $95\%$ interval
$[+51, +237]$, exact permutation $p = 0.039$, ahead in seven of eight).
With the full state observable --- where the task is Markovian and there
is no history to exploit --- the difference is $+24$ and not
distinguishable from zero. The direction of the interaction is as
predicted, though the interaction itself is not resolved at this sample
size.

The return difference is not explained by world-model accuracy. Rolled
forward open-loop on the agent's own trajectories, the two priors
predict equally well (geometric-mean error ratio $1.00$, $p = 1.00$),
and both are roughly three to four times \emph{worse} than a
persistence baseline that predicts no change at all. An agent whose
dynamics model is beaten by a trivial baseline nonetheless gains
$69\%$ in return from changing that model. We trace the decoupling to
the architecture --- the acting agent uses the posterior, so the prior
shapes only the imagined data the policy trains on --- and report it as
a finding about model-based reinforcement learning rather than about
quantum circuits.

We do not claim quantum advantage, and explain quantitatively why none
is available here: a data re-uploading circuit of this depth computes a
truncated Fourier series whose frequency support is fixed by the
encoding, and a classical surrogate spanning that spectrum matches it
with eight parameters against the circuit's thirty. The gains above are
therefore attributable to inductive bias rather than to representational
power, and we state precisely where the bound ceases to apply.

Both circuits run on IonQ Forte-1, where measured parameter-shift
gradients agree with exact simulation. We also quantify a design
constraint that matters for any hardware implementation: at ten qubits
the two standard entangling topologies differ by a factor of
$2{,}766$ in the shot count needed to resolve a gradient. The
architecture is offered as an instrument for the question it makes
askable --- whether quantum dynamics models help when the dynamics are
themselves quantum --- and we release implementation, verification
harness and run records.
\end{abstract}

\vspace{1em}

% =====================================================================
\section{Introduction}
% =====================================================================
Reinforcement learning with quantum components has been almost entirely
model-free. Parameterised quantum circuits (PQCs) have served as
policies and as value-function approximators~\cite{chen2020,skolik2023},
inheriting the sample inefficiency of the model-free setting while
adding the cost of circuit evaluation. The quantum component is placed
where the agent decides, not where it predicts.

Model-based reinforcement learning offers a different insertion point,
and we think a better-motivated one. A model-based agent learns a
\emph{world model} --- a predictive model of environment dynamics ---
and optimises its policy against that model rather than against the
environment. The world model is the component whose job is to represent
dynamics, and dynamics is what a quantum system expresses natively. If
there is a principled place to put a quantum circuit in a reinforcement
learning agent, this is it.

This paper builds that object. We take two constructions studied
separately and join them. From quantum sequence modelling we take the
transformer encoder in which one layer's token mixing is performed by a
PQC~\cite{qasa}. From model-based RL we take the recurrent state-space
model (RSSM) of the DreamerV3 family~\cite{hafner2023}. The join is to
make the quantum transformer the RSSM's \textbf{prior}: the distribution
over the next stochastic latent given the deterministic recurrent state.

The position matters more than it may appear. In latent imagination the
agent rolls the world model forward using the prior alone, because no
future observations exist, and trains actor and critic on the resulting
trajectories. Replacing the prior therefore means that every imagined
transition the policy ever sees during optimisation is generated through
a quantum circuit. The quantum component is not an auxiliary module
consulted occasionally; it is the generative core of the agent's
training signal.

\paragraph{What this paper is for.}
We characterise an instrument. The claims are about what the
construction is, what it predicts, how it behaves on hardware, and what
it costs --- not about beating a classical baseline. That restraint is
deliberate and is grounded in theory rather than in modesty: at the
qubit counts reachable by exact simulation, the function class our
circuit realises is classically reachable, and we demonstrate this
directly with a surrogate (Section~\ref{sec:parity}). A paper reporting
quantum advantage at six qubits on a control task would be reporting an
artefact. What is worth establishing instead is that the architecture
works, that it is correct in the ways that are easy to get silently
wrong, that its predictions are measurably good, and that it transfers
to hardware --- because those are the prerequisites for asking the
question where advantage is not ruled out in advance.

\paragraph{Contributions.}
\begin{enumerate}\itemsep3pt
\item \textbf{A quantum transformer dynamics prior} for recurrent
      state-space models, and the first model-based agent whose world
      model is a quantum transformer, trained end to end through
      multi-step imagination (Section~\ref{sec:arch}).
\item \textbf{A correctness account.} The construction requires a
      bounded history window with exact reset semantics across sequence,
      rollout and episode boundaries; a map from circuit readouts to
      categorical latent logits; gradient flow through $H$ composed
      circuit applications; and attention to KL balancing. Each
      associated failure mode is silent rather than loud, so we specify
      a verification harness that detects them
      (Section~\ref{sec:impl}).
\item \textbf{Predictive performance.} Offline one-step and open-loop
      $k$-step error against a capacity-matched classical twin over
      eight seeds, with an ablation over attention context length at
      fixed model width (Sections~\ref{sec:offline}--\ref{sec:context}).
\item \textbf{A tested architectural prediction.} The attentional
      prior improves return by $144$ points over eight seeds where the
      dynamics are history-dependent ($p = 0.039$) and by an
      indistinguishable $24$ points where they are Markovian, matching
      the pattern the design predicts in advance
      (Section~\ref{sec:closedloop}).
\item \textbf{A sharp decoupling.} Both priors are three to four
      times worse than a persistence baseline as open-loop dynamics
      models, and indistinguishable from each other, while differing by
      $69\%$ in return. Agent performance in this setting is not
      mediated by world-model accuracy
      (Section~\ref{sec:imagination}).
\item \textbf{A quantitative account of why parity is expected},
      including a classical Fourier surrogate that matches the circuit
      with eight parameters, and a statement of the conditions under
      which the argument fails (Section~\ref{sec:parity}).
\item \textbf{Hardware execution} on IonQ Forte-1 with the measurement
      convention settled empirically, plus a $2{,}766\times$ difference
      in gradient-resolution cost between entangling topologies that
      constrains any hardware implementation
      (Section~\ref{sec:hardware}).
\end{enumerate}

% =====================================================================
\section{Background: recurrent state-space models}
% =====================================================================
A model-based agent of the DreamerV3 family~\cite{hafner2023} maintains
a latent state factorised into a deterministic recurrent component
$h_t$ and a stochastic component $z_t$ represented as categorical
variables. Three distributions define the world model: a recurrent
update $h_t = f(h_{t-1}, z_{t-1}, a_{t-1})$; a \emph{posterior}
$q(z_t \mid h_t, o_t)$, which sees the observation; and a \emph{prior}
$p(z_t \mid h_t)$, which does not. Training reconstructs observations
and rewards while keeping the prior close to the posterior under a
balanced Kullback--Leibler term.

\paragraph{Imagination.}
Policy optimisation never touches the environment. The agent samples a
starting latent from replayed experience and rolls the world model
forward for a fixed horizon using \emph{the prior only}. Actor and
critic are trained on these imagined trajectories via
$\lambda$-returns. The prior is thus the generative core of the agent:
it is what the agent dreams with.

\paragraph{Why attention, given a recurrent backbone.}\label{sec:background-attention}
The RSSM already carries history in $h_t$, so it is reasonable to ask
what attention contributes. The answer is that $h_t$ is a \emph{summary}
produced by a gated recurrent unit trained under the world-model loss,
and the baseline prior sees only its instantaneous value. A prior that
attends over $h_{t-K+1}, \dots, h_t$ conditions on the recent
\emph{trajectory} of the recurrent state. This should matter precisely
when the dynamics are history-dependent, which gives the architecture a
falsifiable prediction: the advantage over a non-attentional prior
should appear under partial observation and vanish when the observation
is Markovian. TransDreamer~\cite{transdreamer} makes the same argument
for replacing the recurrent backbone wholesale; we apply it to the prior
alone, so that the comparison against the validated baseline isolates a
single module.

% =====================================================================
\section{Related work}
% =====================================================================
\paragraph{Quantum reinforcement learning.}
Existing work places PQCs in the policy or value
function~\cite{chen2020,skolik2023}. We are not aware of prior work
placing a quantum model in the dynamics model of a model-based agent.
\pend{Appendix~\ref{app:novelty}: databases, search strings, dates}

\paragraph{World models for quantum problems.}
The converse direction exists: model-based RL applied \emph{to} quantum
problems, for instance DreamQAS~\cite{dreamqas}, which learns a
classical world model over circuit-construction trajectories. That work
uses a classical model to design quantum circuits; we use a quantum
circuit inside a model.

\paragraph{Quantum attention.}
QASA~\cite{qasa} introduces the construction our prior adapts. Its
revised version attributes the reported gain to the low-rank bottleneck
the circuit imposes rather than to quantumness as such. We take that
seriously and treat it as an experimental question rather than a
settled one, which is why the bottleneck-width ablation of
Section~\ref{sec:ablation} is part of the design rather than an
afterthought.

\paragraph{Transformers as world models.}
TransDreamer~\cite{transdreamer} replaces Dreamer's recurrent backbone
with a transformer state-space model. We leave the backbone intact, so
every component other than the prior is identical to the validated
baseline.

\paragraph{Expressivity and benchmarking.}
Data re-uploading circuits compute truncated Fourier series in their
inputs, with frequencies fixed by the encoding gates and only the
coefficients trainable~\cite{schuld2021}. This bounds what the circuit
can represent and predicts that a classical model spanning the same
frequency support will match it; Section~\ref{sec:parity} confirms this
for our circuit. Bowles et al.~\cite{bowles2024} find that quantum
models at simulable scale generally match rather than exceed classical
baselines of comparable capacity, and that removing entanglement often
leaves performance unchanged. Both findings inform our choice of
controls.

% =====================================================================
\section{Architecture}\label{sec:arch}
% =====================================================================
\subsection{The quantum transformer prior}
In the baseline RSSM the prior is a multilayer perceptron from the
deterministic recurrent state to logits over the categorical stochastic
latent,
\begin{equation}
  \mathrm{logits} = \mathrm{MLP}(h_t) \in \mathbb{R}^{S \times C},
  \qquad
  p(z_t \mid h_t) = \mathrm{Categorical}(\mathrm{logits}),
\end{equation}
for $S$ stochastic variables with $C$ classes each. We replace this map
with a quantum transformer over a bounded window of recent recurrent
states. Writing $H_t = (h_{t-K+1}, \dots, h_t)$,
\begin{equation}
  \mathrm{logits} = W_{\mathrm{out}}\,
  \Big[\mathrm{Enc}_{\mathrm{q}}\!\big(
      \mathrm{LN}(W_{\mathrm{in}} H_t) + \mathrm{PE}\big)\Big]_{-1},
\end{equation}
where $\mathrm{PE}$ is sinusoidal positional encoding, the final token's
representation is read out, and $\mathrm{Enc}_{\mathrm{q}}$ is a
two-block transformer encoder in which block~$1$ carries the PQC. The
interface $h_t \mapsto \mathrm{logits}$ is preserved exactly, so the
recurrent update, posterior, decoders, actor, critic, KL balancing and
imagination loop are untouched.

Each encoder block applies causally masked multi-head self-attention,
then --- in the quantum block --- the circuit as a residual token-mixing
sublayer, then a position-wise feed-forward network, each with layer
normalisation:
\begin{align}
  u &= \mathrm{LN}_1\big(x + \mathrm{MHA}(x, x, x)\big), \\
  v &= \mathrm{LN}_2\big(u + \mathrm{QuantumToken}(u)\big), \\
  y &= \mathrm{LN}_3\big(v + \mathrm{FFN}(v)\big).
\end{align}
$\mathrm{QuantumToken}$ acts on each sequence position independently:
project $d_{\mathrm{model}} \to n$ rotation angles by a linear map
followed by $\tanh$ and scaling by $\pi$; evaluate the circuit; project
the $n$ readouts back to $d_{\mathrm{model}}$.

\paragraph{The circuit.}
Angle encoding with data re-uploading over $L = 4$ layers on
$n = 6$ data qubits plus one ancilla, brickwork entangler, local
Pauli-$Z$ readout giving $n$ outputs in $[-1, 1]$. Trainable circuit
parameters number $L \times (n+1) = 28$. The entangler and readout
choices are not stylistic: with a brickwork pattern and local
observables, measured gradient variance is flat in width out to ten
qubits, whereas a sequential CNOT ring with global readout decays by a
factor of roughly $0.5$ per added qubit~\cite{papera}. Brickwork is
genuinely depth-two at any width, so the causal cone grows with layer
count as intended~\cite{cerezo2021}.

\begin{table}[htb]\centering\small
\caption{The dynamics prior and its capacity-matched classical twin.
The two arms differ in exactly one sublayer.}
\label{tab:arch}
\begin{tabular}{lll}
\toprule
& Quantum transformer prior & Classical twin \\
\midrule
Attention context $K$          & 4                  & 4 \\
$d_{\mathrm{model}}$           & 64                 & 64 \\
Attention heads                & 4                  & 4 \\
Encoder blocks                 & 2                  & 2 \\
Token-mixing block index       & 1                  & 1 \\
Token-mixing map               & PQC, 6 qubits      & $\tanh$ MLP, width 6 \\
Data qubits $+$ ancilla        & $6 + 1$            & --- \\
Circuit layers $L$             & 4                  & --- \\
Encoding                       & angle, re-uploaded & --- \\
Entangler                      & brickwork          & --- \\
Readout                        & local Pauli-$Z$    & --- \\
\midrule
Trainable circuit parameters   & 28                 & 0 \\
Total prior parameters         & 84{,}770           & 85{,}494 \\
\bottomrule
\end{tabular}
\end{table}

\subsection{The history window}
The prior is called once per RSSM step by code that knows nothing about
attention, so the window is maintained as state on the module and
advanced on each call. This is the component where an error would be
invisible. A window not cleared at an episode boundary would let one
episode attend over the tail of the previous one; a window not cleared
between imagination rollouts would let imagined futures attend over
replayed history. Either would inflate apparent performance without
raising an error.

We therefore install reset logic by wrapping the three entry points at
which a boundary occurs --- sequence start in \texttt{observe},
per-step reset in \texttt{obs\_step}, and rollout start in
\texttt{imagine} --- on the agent instance, leaving the baseline
implementation unmodified. The semantics are asserted by test
(Section~\ref{sec:impl}), including the two cases that are easy to miss:
a batch-size change between the observation and imagination phases,
which must re-seed rather than reuse the window, and a partial reset in
which only some rows of the batch have reached a boundary.

\subsection{What the circuit is load-bearing for}
The circuit is $28$ of the prior's $84{,}770$ parameters. The claim the
architecture supports is therefore one of \emph{information bottleneck},
not of parameter share: in the quantum block, every token's
representation is routed through $n = 6$ real numbers produced by a
quantum circuit, on the residual path of the only block that mixes
tokens that way, and the surrounding layers cannot route around it.

Whether the circuit or the surrounding width carries the predictive work
is an empirical question with a decisive experiment, and we run it
rather than assume the favourable answer
(Section~\ref{sec:ablation}).

\subsection{Controls}
\emph{Capacity-matched classical twin.} The identical transformer with
the PQC sublayer replaced by a $\tanh$ MLP of the same bottleneck width,
in the same position in the same block. Attention, positional encoding,
layer norms, feed-forward network, window semantics and context length
are shared code. The arms match to within $1\%$ on parameter count
($84{,}770$ against $85{,}494$), and the test suite fails if they
diverge by more than $15\%$.

\emph{Non-attentional prior.} The construction without attention or
history: down-projection, circuit, expansion head. Retained as the
ablation that isolates what the attention contributes, with
$18{,}658$ parameters against the twin's $18{,}672$. Both replace
DreamerV3's original $131{,}840$-parameter prior, so both arms are
narrower at that position than the unmodified baseline --- a property of
the comparison worth keeping in view.

\emph{Entanglement removed.} All two-qubit gates deleted, parameter
count unchanged, so the comparison is otherwise exact. Following
Bowles et al.~\cite{bowles2024}, an architecture paper should establish
this of its own construction.

% =====================================================================
\section{Implementation and verification}\label{sec:impl}
% =====================================================================
The failure modes in this integration are silent. A missing gradient
path, an untracked parameter group or a leaking history window produces
runs that complete normally and report plausible numbers. We therefore
treat verification as part of the contribution.

\paragraph{Optimiser reconstruction.}
The prior is installed by substituting a module into an
already-constructed agent. An optimiser built before the substitution
holds references to the \emph{replaced} parameters, so the circuit
parameters are silently untracked and training proceeds without error
while the quantum component never updates. The optimiser must be rebuilt
after substitution.

\paragraph{Gradient flow through imagination.}
The actor is trained on trajectories produced by rolling the prior
forward for $H$ steps, so gradients reach the circuit through a
composition of $H$ applications of the same circuit. We instrument
gradient magnitudes at the circuit parameters separately from the
surrounding classical ones, extending the barren-plateau baseline
through training rather than measuring it once at initialisation.

\paragraph{KL balancing.}
The world-model loss regularises the prior toward the posterior with
asymmetric weighting. A prior with a materially different
parameterisation changes the effective strength of this term, so a
systematic difference between the arms would confound the return
comparison of Section~\ref{sec:closedloop}. We therefore report it.

\figorplaceholder{fig_kl.pdf}{0.62}{Prior--posterior Kullback--Leibler
divergence through training, partial observation, median across seeds
with interquartile band. This is the quantity the world-model loss
regularises. Comparable trajectories indicate the two arms are being
regularised to a similar degree, so the return difference is not an
artefact of one prior sitting at a different point in the KL trade-off.}{kl}

\paragraph{Verification harness.}
Nine checks must pass before any training run, for both arms:
output shape and finiteness; the window is bounded at $K$; history
changes the prediction, so the attention is not decorative;
\texttt{reset()} fully clears, so two fresh sequences agree exactly;
residual nondeterminism in training mode is attributable to attention
dropout and not to retained state; selective row reset is complete and
leaves other rows bit-identical; a batch-size change re-seeds the
window; gradients reach the circuit parameters specifically with
non-zero magnitude; and the two arms are capacity-matched.

Two of these require care. The history and reset checks must run under
evaluation mode: with dropout active, two identical calls differ
regardless of whether the window was cleared, so the history check would
pass on noise while the reset check would fail on a correct
implementation. We verify separately that the module is stochastic under
training mode and deterministic under evaluation mode, which localises
the nondeterminism to dropout.

\paragraph{Run provenance.}
Every configuration setting appears unconditionally in the run directory
name, training refuses to write into a directory containing a completed
run, and the sweep driver refuses to start when a result-affecting
setting absent from the directory name differs from its default. These
guards exist because their absence caused two silent data losses during
this work; Appendix~\ref{app:notes} records what happened, on the view
that a reader cannot otherwise rule such failures out.

% =====================================================================
\section{Predictive performance}\label{sec:offline}
% =====================================================================
The prior's function is to predict latent dynamics. We measure that
directly, in isolation, before asking what it does inside an agent.
This is also where the statistical resolution is: an offline seed costs
roughly half an hour against twenty-seven hours for a closed-loop seed,
so eight seeds and a context ablation are affordable
(Section~\ref{sec:power}).

\paragraph{Setup.}
Latent trajectories are collected from a trained DreamerV3 agent on
CartPole under a mixed policy ($70\%$ greedy, $30\%$
$\epsilon$-random), which gives the state coverage a dynamics model
needs; a purely greedy policy yields a narrow on-policy manifold that
flatters every model equally. Dreamer's $[h_t \mid z_t]$ features are
compressed to six dimensions by an autoencoder --- one feature per
qubit under angle encoding --- with held-out explained variance of
approximately $91\%$. Splits are by episode rather than by window, since
windowing across a split boundary leaks future frames into training.

Two metrics: one-step prediction error, and $k$-step \emph{open-loop}
rollout error in which the model re-consumes its own predictions. The
latter is the regime a world model is actually used in; teacher-forced
error looks far better and is not what the imagination loop does. The
pre-registered gate is to match the classical twin at one step and stay
within $2\times$ at ten steps.

\begin{table}[htb]\centering\small
\caption{Offline prediction on held-out CartPole episodes,
$d_{\mathrm{model}} = 32$, eight seeds per cell, paired by seed. Ratios
are quantum/classical; below $1$ favours the quantum prior. $p$ is an
exact sign-flip permutation test on the paired differences.}
\label{tab:offline}
\begin{tabular}{lcccccc}
\toprule
& \multicolumn{3}{c}{1-step} & \multicolumn{3}{c}{10-step open-loop} \\
\cmidrule(lr){2-4}\cmidrule(l){5-7}
Context $K$ & ratio & $p$ & Q ahead & ratio & $p$ & Q ahead \\
\midrule
$2$  & $0.984$ & $0.224$ & $4/8$ & $1.007$ & $0.860$ & $2/8$ \\
$4$  & $0.988$ & $0.290$ & $5/8$ & $0.992$ & $0.907$ & $3/8$ \\
$8$  & $0.975$ & $0.041$ & $6/8$ & $0.908$ & $0.283$ & $4/8$ \\
$16$ & $1.014$ & $0.091$ & $2/8$ & $1.107$ & $0.103$ & $2/8$ \\
\bottomrule
\end{tabular}
\end{table}

\paragraph{The two priors are equally good offline.}
Ratios span $0.975$ to $1.014$ at one step and $0.908$ to $1.107$ at ten.
The largest effect anywhere in the table is $2.5\%$. One cell reaches
$p = 0.041$, which is one of eight tests: the Bonferroni threshold is
$0.0063$, and under the global null one would expect $0.4$ cells below
$0.05$. We therefore read Table~\ref{tab:offline} as parity, and do not
claim the $K = 8$ cell as a finding.

At $d_{\mathrm{model}} = 6$ the quantum arm is slightly worse
($1.031\times$ at one step, ahead in two of eight), so the parity does
not depend on choosing a favourable width.

This is the outcome Section~\ref{sec:parity} predicts. At six qubits the
circuit's accessible function class is classically reachable, so a
capacity-matched classical model should match it, and it does.

% =====================================================================
\section{Attention context length}\label{sec:context}
% =====================================================================
Section~\ref{sec:background-attention} gave the architecture a
falsifiable prediction: a prior that attends over recent recurrent
states should benefit from a longer window in proportion to how much
history the dynamics require. Table~\ref{tab:offline} tests it directly
at fixed width, and \textbf{does not support it}.

\paragraph{No trend.}
The one-step ratio moves $0.984 \to 0.988 \to 0.975 \to 1.014$ across
$K = 2, 4, 8, 16$, and the ten-step ratio
$1.007 \to 0.992 \to 0.908 \to 1.107$. Neither is monotonic, and the
longest window is the worst cell in both rows.

\paragraph{Context helps --- but it helps the classical twin more.}
Absolute ten-step error improves with $K$ in both arms, which says the
task genuinely rewards history. Between $K = 2$ and $K = 16$ the quantum
arm improves from $0.00564$ to $0.00527$, a gain of $7\%$; the classical
arm improves from $0.00560$ to $0.00476$, a gain of $15\%$. A trend
present in both arms is a property of the task rather than of the
circuit, and here the arm \emph{without} the circuit extracts more from
the additional context. We report this because it is the result a
referee would otherwise derive from our own table.

\figorplaceholder{context_cartpole_d32.png}{0.62}{Open-loop ten-step
error against attention context $K$ at fixed width
($d_{\mathrm{model}} = 32$), eight seeds per point. Both arms improve
with a longer window, which indicates the task genuinely rewards
history; the classical twin improves more. Generated by the Phase~1
aggregation script.}{context}

\paragraph{What we conclude.}
Attention over recent recurrent states is useful --- both arms improve
with it --- but there is no evidence that routing the token mixing
through a quantum circuit makes it more so. The architectural
hypothesis, as a claim about \emph{offline predictive accuracy}, is not
supported by eight seeds at four context lengths.

This matters for how Section~\ref{sec:closedloop} should be read. The
closed-loop difference reported there cannot be explained by the quantum
prior being a more accurate dynamics model, because on identical
held-out data it is not one. Section~\ref{sec:imagination} takes up what
that leaves.

% =====================================================================
\section{Closed-loop behaviour}\label{sec:closedloop}
% =====================================================================
\subsection{Setup}
CartPole, with a training configuration reduced relative to DreamerV3's
defaults: batch $8$ against $16$, sequence length $16$ against $32$,
imagination horizon $10$ against $15$, and a gradient step every four
environment steps rather than every two. The reduction is forced by
circuit cost --- imagination performs \texttt{horizon} sequential
circuit calls over $\texttt{batch} \times \texttt{seq\_len}$ states per
gradient step --- and applies identically to both arms. The solved
threshold is a mean evaluation return of $475$.

Final returns are measured over $50$ evaluation episodes from the saved
checkpoint rather than the $20$-episode figure recorded during training;
at $20$ episodes the standard error alone moves individual seeds by up
to $90$ points.

\paragraph{Reference implementation.}
Full DreamerV3 at its own default settings solves CartPole in all three
seeds we ran, first crossing the threshold at $10{,}000$, $15{,}000$ and
$45{,}000$ environment steps. Our reduced configuration therefore costs
a factor of roughly two to eight in sample efficiency, which is the
context for the step budget below.

\paragraph{Why $K = 4$.}
Measured throughput is $2.8$ environment steps per second at $K = 4$
against $0.4$ at $K = 8$ --- a factor of seven, and the difference
between an eight-hour run and a fifty-five-hour one. With a fixed
compute budget, four seeds at $K = 8$ buy less resolution than
sixteen at $K = 4$, and Section~\ref{sec:power} shows resolution is the
binding constraint.

\subsection{The prediction the architecture makes}
Full-state CartPole is Markovian: the observation determines the next
state, so a history-dependent prior has little to contribute. Hiding
both velocity components leaves only cart position and pole angle, which
makes the latent dynamics genuinely history-dependent while changing
nothing else about the task, the budget or the architecture.

Section~\ref{sec:context} stated the resulting prediction before these
runs were analysed: the quantum transformer prior should be ahead under
partial observation and at parity under full observation. Both
conditions use the same architecture, the same context length and the
same capacity-matched control, so this is a directional test of a
stated hypothesis rather than a comparison assembled after the fact.

\subsection{Results}
\begin{table}[htb]\centering\small
\caption{Final return at $80{,}000$ environment steps, $50$ evaluation
episodes per agent, quantum transformer prior against its
capacity-matched classical twin. Both observability conditions use
$K = 4$.}
\label{tab:closedloop}
\begin{tabular}{lrrrrrrrrr}
\toprule
& \multicolumn{9}{c}{Seed} \\
\cmidrule(l){2-10}
& 0 & 1 & 2 & 3 & 4 & 5 & 6 & 7 & Mean \\
\midrule
\multicolumn{10}{l}{\emph{Partial observation} (velocities hidden)} \\
Quantum    & $420.9$ & $375.2$ & $478.7$ & $313.5$ & $368.9$ & $298.9$ & $131.5$ & $443.9$ & $\mathbf{353.9}$ \\
Classical  & $111.0$ & $312.6$ & $469.7$ & $195.0$ & $150.8$ & $138.0$ & $205.0$ & $96.9$  & $209.9$ \\
Difference & $+309.9$ & $+62.6$ & $+9.0$ & $+118.5$ & $+218.1$ & $+160.9$ & $-73.5$ & $+347.0$ & $\mathbf{+144.1}$ \\
\midrule
\multicolumn{10}{l}{\emph{Full observation} (Markovian)} \\
Quantum    & $495.6$ & $459.1$ & $196.0$ & $396.1$ & $205.3$ & $233.6$ & --- & --- & $330.9$ \\
Classical  & $228.5$ & $391.9$ & $139.3$ & $469.8$ & $184.8$ & $426.2$ & --- & --- & $306.8$ \\
Difference & $+267.1$ & $+67.2$ & $+56.7$ & $-73.7$ & $+20.5$ & $-192.6$ & --- & --- & $+24.2$ \\
\bottomrule
\end{tabular}
\end{table}

\paragraph{Partial observation.}
The quantum transformer prior leads by $144.1$ return points, a $69\%
$ relative improvement ($353.9$ against $209.9$), ahead in seven of
eight seeds. The $95\%$ bootstrap interval is $[+51.0, +236.8]$ and
excludes zero; an exact sign-flip permutation test gives
$p = 0.039$ against an attainable floor of $0.008$ at this sample size.
From the observed dispersion (standard deviation of the paired
difference $145.1$) the seeds required for $80\%$ power are
approximately eight, so the comparison is powered at roughly the
level it was run --- marginally, but not retrospectively.

\paragraph{Full observation.}
The difference is $+24.2$ with interval $[-86.2, +135.9]$ and
$p = 0.656$, ahead in four of six seeds. This is parity. It is also what
Section~\ref{sec:parity} and the Markovian structure of the task
jointly predict: with the full state observable there is no history for
an attentional prior to exploit, and no representational advantage
available at this circuit scale.

\paragraph{The interaction.}
The swing between conditions is $+119.9$ return points
($+144.1$ against $+24.2$), with a standard error of $81.1$ and a
two-sample permutation $p$ of $0.162$ on the difference of differences.
The direction is as predicted and the main effect under partial
observation is resolved, but \textbf{the interaction itself is not}.
Interaction tests need roughly four times the sample of a main effect,
and six and eight seeds do not supply it. We therefore claim the
partial-observation result and report the crossover as consistent with
the architectural hypothesis rather than as established by it.

\paragraph{Two cautions on the partial-observation result.}
First, multiplicity. This comparison is one of four we report
(two priors $\times$ two observability conditions). Under a Bonferroni
correction for that family, $p = 0.039$ would not survive
$\alpha = 0.0125$. Our defence is that the comparison is directional and
was specified from the architecture in advance rather than selected
after inspection --- but a reader should weigh the claim accordingly,
and we note that the same hypothesis tested on the non-attentional prior
returned null (Section~\ref{sec:power}).

Second, the dispersion is still large in absolute terms: per-seed
differences range from $-73.5$ to $+347.0$. The effect is resolved
because it is large, not because the measurement is quiet. Four
additional seeds per arm would move the comparison from marginally to
comfortably powered, and we regard that as the first extension worth
making.
\figorplaceholder{fig_learning.pdf}{1.0}{Evaluation return against
environment steps, median across seeds with interquartile band;
individual seeds faint behind. Left: partial observation, where the
quantum transformer prior leads. Right: full observation, where the two
arms are indistinguishable. Dotted line marks the solved threshold of
$475$, which neither arm reaches under the reduced training
configuration of Section~\ref{sec:closedloop}. We plot the median rather
than the mean because the per-seed distribution is strongly skewed --- at
partial observation one seed finishes at $131.5$ and another at
$478.7$ --- so a mean curve would imply a symmetry the data does not
have.}{learning}

\figorplaceholder{fig_gradient.pdf}{0.62}{Variance of the gradient on
the $28$ circuit parameters through training, both observability
conditions. This extends the barren-plateau diagnostic across training
rather than measuring it once at initialisation. A curve decaying
steadily toward zero would indicate that the circuit had stopped
learning while the surrounding classical layers continued to improve ---
the failure mode that would make a nominally quantum result actually
classical.}{gradient}
\pend{confirm the gradient instrumentation fired during these runs; if
the \texttt{prior\_q\_var} column is empty this figure must be withdrawn
rather than reported}

% =====================================================================
\section{World-model quality, and a sharp decoupling}
\label{sec:imagination}
An agent can reach competence while effectively ignoring a poor world
model, learning a policy that works despite its predictions rather than
because of them. We therefore measure the prior directly: roll it
forward without observations, as imagination does, and compare the
decoded predictions against the observations actually received. All
figures are in $\mathrm{symlog}$ space, where the DreamerV3 decoder is
trained. We add a \emph{persistence} baseline --- predict no change ---
because its omission would flatter any learned model.

\paragraph{Scope of the window.}
Each episode is truncated at $22$ steps from reset, so what follows
measures short-horizon prediction during the recovery phase rather than
over a full episode. This is a deliberate restriction: once a policy has
stabilised, ``no change'' is close to correct and would dominate the
comparison for reasons that say nothing about either model. The cost is
that these numbers speak only to the transient regime. Extending the
measurement over full episodes, with persistence reported alongside, is
straightforward and not done here.

\begin{table}[htb]\centering\small
\caption{Open-loop prediction error at horizon $10$ on the agent's own
trajectories, partial observation, eight seeds. ``Ratio to
persistence'' above $1$ means the learned prior predicts worse than
assuming nothing changes.}
\label{tab:imagination}
\begin{tabular}{lccc}
\toprule
& Quantum prior & Classical twin & Quantum/classical \\
\midrule
Mean $k{=}10$ error      & $0.00426$ & $0.00468$ & $1.00$ \\
Geometric-mean ratio     & \multicolumn{2}{c}{---} & $1.00$, $p = 1.00$ \\
Quantum better in        & \multicolumn{2}{c}{---} & $4/8$ seeds \\
\midrule
Median ratio to persistence & $3.81$ & $3.14$ & \\
Worse than persistence in   & $8/8$  & $8/8$  & \\
\bottomrule
\end{tabular}
\end{table}

Two results, and neither is the one we expected.

\paragraph{The two priors are equally good dynamics models.}
The geometric-mean quantum/classical error ratio is $1.00$ with an exact
permutation $p$ of $1.00$ and the quantum arm ahead in four of eight
seeds. Under full observation the ratio is $1.08$ ($p = 0.84$, ahead in
two of six). Whatever produces the $144$-point return difference of
Section~\ref{sec:closedloop}, it is not a difference in one-shot
open-loop predictive accuracy as measured here.

\paragraph{Both are worse than predicting no change.}
Every one of the sixteen partial-observation runs has a
persistence ratio above $1$, with medians near $3$--$4$.

The obvious objection is that persistence is artificially strong because
a competent agent holds the pole still, so ``nothing changes'' is nearly
correct. That objection does not apply here, for two independent
reasons. First, the measurement truncates each episode at $22$ steps
from reset, so it observes only the recovery phase in which the state is
actively moving, and never the stabilised phase of a $400$-step
episode --- the regime in which persistence would be strong is excluded
by construction. Second, and consistently with that, the rank
correlation between final return and persistence ratio across the
sixteen runs is $+0.13$ (Pearson $-0.02$): the ratio does not track
agent competence.

Persistence is therefore a \emph{weak} baseline on these trajectories,
and both learned priors lose to it regardless. We report this without a
demonstrated cause.

A plausible mechanism, which we flag rather than claim: the prior is
trained by a one-step Kullback--Leibler term against the posterior and
is never optimised for ten-step open-loop rollout, so compounding error
in latent space is unconstrained by the training objective. This is a
property of the DreamerV3 objective rather than of the quantum
substitution, consistent with both arms behaving the same way.

\paragraph{What follows: the decoupling is the finding.}
Taken together these results sharpen rather than weaken the paper's
central empirical observation. An agent whose dynamics model is
three to four times worse than a trivial baseline still gains $69\%$ in
return when that model's token-mixing sublayer is changed, and the two
models are indistinguishable on the accuracy metric. Agent performance
here is not mediated by open-loop world-model accuracy.

The mechanism is architectural and not specific to quantum circuits:
when acting, the agent uses the \emph{posterior}, which sees real
observations. The prior shapes only the imagined trajectories on which
actor and critic are trained. What matters for the policy is evidently
not how accurately those trajectories predict the future, but some
other property of them --- plausibly their diversity, their smoothness,
or the shape of the value landscape they induce. Identifying which is
beyond what our data can settle, and we name it as an open question.

Two consequences for method. First, a quantum substitution inside a
world model cannot be assessed by open-loop accuracy alone; the
measurement is nearly blind to a change that moves return by $69\%$.
Second, and conversely, it cannot be assessed by return alone either,
at the dispersion documented in Section~\ref{sec:power}. Both
measurements are necessary and neither is sufficient, which is an
uncomfortable methodological position and, we think, the correct one.

\paragraph{Where the clean dynamics comparison lives.}
The measurement above is taken on each agent's \emph{own} final-policy
trajectories, which differ between arms and are narrow by construction
once a policy has stabilised. It is the right measurement for asking
``does this agent use its model'', and the wrong one for asking ``which
prior is the better dynamics model''. The latter question is answered
by the offline protocol of Section~\ref{sec:offline}, where both arms
see identical held-out episodes drawn from a mixed exploratory policy,
and where eight seeds are affordable.

\section{Why parity is the expected outcome}\label{sec:parity}
% =====================================================================
Closed-loop returns at parity with a capacity-matched classical control
should not be read as a shortfall of the architecture. It is what
expressivity theory requires at this scale, and the argument is specific
enough to be checked.

\paragraph{The bound.}
A data re-uploading circuit computes a truncated Fourier series in its
inputs~\cite{schuld2021}. The accessible frequency spectrum is fixed
entirely by the encoding gates and their repetition --- for our circuit,
angle encoding re-uploaded over $L$ layers gives integer frequencies up
to $L$ per input dimension --- while the trainable parameters set only
the coefficients. The circuit is therefore a linear model over a fixed,
known, and modestly sized feature basis.

\paragraph{The check.}
If that account is right, a classical model spanning the same spectrum
should match the circuit. We built one: a layer with fixed random
integer frequencies drawn from the circuit's accessible range and
trainable coefficients, substituted for the PQC with everything else
held constant. It matches the circuit's error using \textbf{eight}
trainable parameters against the circuit's \textbf{thirty}. Sweeping the
surrogate's capacity over a $64\times$ range moves its error less than
changing the random seed does~\cite{papera}.

\paragraph{What follows.}
Two things. First, at six qubits and depth four, any performance
difference between our quantum prior and a classical model of matched
capacity must come from optimisation dynamics, inductive bias or
regularisation --- not from representational power, because the
representation is classically reachable by construction. Reporting
quantum advantage in this regime would be reporting an artefact, and
parity is the correct and informative result. Second, and more usefully,
the argument identifies exactly what would have to change for it to
stop applying:

\begin{itemize}\itemsep3pt
\item \textbf{Qubit count beyond exact simulation.} The surrogate is
      constructible because we can enumerate the spectrum. That
      enumeration grows with the encoding, not with the Hilbert space,
      so the relevant threshold is where the spectrum becomes
      intractable rather than where the statevector does.
\item \textbf{Dynamics that are themselves quantum.} The Fourier basis
      is efficient for the smooth, low-frequency dynamics of a classical
      control task. For a target function that is natively expressed in
      the circuit's own structure --- the time evolution of a spin
      system, say --- the classical surrogate must approximate what the
      circuit computes directly. This is the condition under which the
      bound does not rule advantage out, and it is where we take the
      architecture next.
\item \textbf{Encodings with richer spectra.} Amplitude encoding and
      trainable encodings change the accessible frequencies rather than
      the coefficients, and are not covered by the surrogate we built.
\end{itemize}

Stating the bound precisely is more valuable than evading it. It tells a
reader what this architecture cannot do at this scale, and it converts
"try a bigger circuit" into a specific experiment.

% =====================================================================
\section{Resolving power}\label{sec:power}
% =====================================================================
Closed-loop returns in this setting carry large seed-to-seed dispersion,
which sets a hard limit on what any study of this size can conclude from
them. We quantify it rather than work around it, because the limit
determines where the paper's quantitative weight should sit.

For the non-attentional prior under partial observation, the paired
difference and its dispersion evolve as seeds accumulate: $+212.2$ at
three seeds (standard deviation $117$), $+122.1$ at six ($182$), and
$+58.3$ at seven ($237$). Dispersion grows faster than the effect, and
the seeds required to detect the then-current effect at $80\%$ power
move from $3$ to $18$ to roughly $130$. At twenty-seven hours per run
that last figure is months of continuous computation for one cell of one
table.

Two consequences follow, and both are design decisions rather than
disclaimers. First, the architecture's quantitative claims belong in the
offline measurements, where a seed is fifty times cheaper and eight
seeds are routine. Second, small-seed closed-loop comparisons must be
reported with their attainable resolution: at three seeds the smallest
two-sided $p$-value an exact sign-flip permutation test can return is
$0.250$, so such a study cannot reach significance whatever the data ---
a fact more useful to a reader than any $p$-value computed from it.

This is not peculiar to our setup. A companion study~\cite{papera}
measures seed dispersion across five quantum sequence-modelling tasks
and finds the sample size required for a $10\%$ effect spanning $6$ to
$150$ depending on task; that under the paired null, $29\%$ of
five-seed studies report a spurious gain of $10\%$ or more; and that a
three-seed pilot cannot forecast its own required sample size, the same
pilot data being consistent with anywhere between $4$ and $170$ seeds.

% =====================================================================
\section{Execution on quantum hardware}\label{sec:hardware}
% =====================================================================
Both circuits were executed on IonQ Forte-1 at eight data qubits (nine
wires), two layers and $504$ shots, using the parameter-shift rule: two
circuit evaluations per gradient, four jobs in total.

\begin{table}[htb]\centering\small
\caption{Parameter-shift gradients measured on IonQ Forte-1 against
exact simulation.}
\begin{tabular}{lccc}
\toprule
Entangler & Hardware & Exact & Difference \\
\midrule
Brickwork    & $-0.6230$ & $-0.5533$ & $-0.0697$ \\
Ring cascade & $+0.0734$ & $+0.0380$ & $+0.0354$ \\
\bottomrule
\end{tabular}
\end{table}

Both agree with exact simulation within measurement error: the
construction transfers to hardware. Our reconstruction of
$\langle Z_0 \rangle$ from the returned probability histograms
reproduces the provider plugin's own expectation values exactly, as
rational multiples of $1/504$, which settles the bit-ordering convention
empirically rather than by assumption --- worth stating because an
inverted convention would produce a plausible disagreement attributable
to noise.

The measured root-mean-square deviation from exact over the four
evaluations is $0.0594$, against $0.0345$ predicted by the provider's
own Forte-1 noise model scaled to this shot count. We report this as
indicative only: it is a single statistic over four evaluations at one
parameter point, carrying roughly $35\%$ relative uncertainty, and the
chosen point has a ring-cascade gradient about $2.5\times$ larger than
the landscape average. The hardware run establishes agreement with
simulation; it is not evidence about typical gradient magnitudes.

\paragraph{A design constraint for hardware implementations.}
The two standard entangling topologies differ enormously in how
expensive their gradients are to resolve. Because shot noise on an
estimated expectation value scales as $1/\sqrt{N}$, the shots needed to
resolve a gradient of magnitude $g$ scale as $g^{-2}$, and gradient
magnitude depends strongly on the entangler. At ten qubits with local
readout, resolving a gradient to fixed relative precision requires
approximately $3$ shots with a brickwork entangler and $9{,}360$ with a
sequential CNOT ring --- a factor of $2{,}766$~\cite{papera}. The ring
is the more common choice in the literature. For simulation this is a
stylistic difference; for hardware it is the difference between a
feasible experiment and an impossible one, and it is the single most
consequential design decision in the circuit.

\paragraph{Scope of the hardware experiment.}
Four jobs, at a total cost of \$672.80. Training on hardware is not
merely expensive but infeasible: at the parameter and step counts of one
$80{,}000$-step run, the parameter-shift rule would require on the order
of $10^{12}$ shots. Hardware execution here is inference-only, and is
intended to establish that the circuits transfer rather than that the
training does.

% =====================================================================
\section{Design space}\label{sec:ablation}
% =====================================================================
\begin{pendingbox}{bottleneck width and entanglement removal}
Both sweeps have been run on the non-attentional prior at $40{,}000$
steps and two seeds. Two seeds is below the resolution
Section~\ref{sec:power} shows is interpretable, so these are reported as
directional pending extension or restatement offline, where seeds are
cheap.

\textbf{Bottleneck width.} The width surrounding the circuit, reducing
the expansion head from $17{,}088$ parameters to $1{,}308$. Performance
roughly flat as it shrinks shows the circuit carries the predictive
work; a collapse shows the surrounding width does, and would require the
bottleneck claim of Section~\ref{sec:arch} to be restated. Given that
QASA's revised version attributes its own gain to a bottleneck rather
than to quantumness~\cite{qasa}, this is the decisive ablation.

\textbf{Entanglement removed.} All two-qubit gates deleted, parameter
count unchanged. Bowles et al.~\cite{bowles2024} report that removing
entanglement frequently leaves performance unchanged; we establish this
for our own construction.

\textbf{Quantum block position.} We adopt QASA's finding that position
matters more than count without re-testing it in this setting, and flag
it as inherited rather than established here.
\end{pendingbox}

% =====================================================================
\section{Discussion}
% =====================================================================
\paragraph{What the construction establishes.}
A quantum transformer can serve as the generative dynamics model of a
model-based agent, trained end to end through multi-step imagination,
with the policy learning entirely from trajectories generated through a
quantum circuit. The construction is stable, reaches competence, admits
a capacity-matched classical control differing in exactly one sublayer,
and transfers to trapped-ion hardware. Measured as a dynamics model
under partial observation it is indistinguishable from that control,
yet as an agent in the same condition it returns $69\%$ more reward
over eight seeds.

That last result deserves a precise reading. It is not evidence of
quantum advantage, for the reason given in
Section~\ref{sec:parity}: the circuit's function class is classically
reachable at this scale, so a classical model of matched capacity could
in principle realise the same mapping. What the comparison isolates is
therefore an \emph{inductive bias} --- the two arms differ in one
sublayer and in nothing else, and the arm containing the circuit
generalises better from the same data under the same budget. Inductive
bias is a weaker claim than advantage and a more useful one, because it
is the property that would survive into regimes where the
representational argument no longer applies.

\paragraph{What it bounds.}
At six qubits and depth four the circuit's function class is classically
reachable, and we demonstrate this with an eight-parameter surrogate.
Parity in closed-loop return is therefore the predicted outcome rather
than a disappointment, and the paper states the three conditions ---
spectrum intractability, natively quantum dynamics, richer encodings ---
under which the bound ceases to apply.

\paragraph{What it enables.}
The value of an instrument is the question it makes askable. The
question here is whether a quantum dynamics model helps when the
dynamics are \emph{themselves} quantum, which is precisely the regime
the Fourier argument does not cover. That experiment requires exactly
what this paper builds: a dynamics prior with a quantum circuit in it, a
capacity-matched classical control, verified correctness, and a measured
account of how many seeds a conclusion needs. We regard the architecture
as the prerequisite rather than the result.

A second contribution is methodological and transferable: the
decoupling of Section~\ref{sec:imagination} implies that quantum
substitutions inside a world model should be evaluated by measuring the
world model. Studies that substitute a component and report only agent
return are measuring through a low-pass filter, with the resolution cost
quantified in Section~\ref{sec:power}.

% =====================================================================
\section{Limitations}
% =====================================================================
\begin{itemize}\itemsep3pt
\item One architecture family, six qubits, simulated exactly during
      training. Hardware execution was inference-only.
\item The partial-observation result is powered at approximately the
      level it was run (eight seeds against a requirement of eight) and
      would not survive a Bonferroni correction across the four
      comparisons reported. It is a directional test of a stated
      hypothesis, not a screen, but four further seeds per arm would
      settle it.
\item The crossover interaction is directionally as predicted but
      unresolved ($p = 0.162$). Interactions require roughly four times
      the sample of a main effect.
\item One environment, low-dimensional control. Scaling to pixel
      observations, where world models matter most, is untested. A
      second environment exists in the codebase but its
      continuous-action paths are unverified.
\item The bottleneck-width and entanglement ablations are at two seeds
      and are reported as directional.
\item The quantum prior costs roughly $12\times$ the classical twin in
      wall clock on the offline task. \pend{closed-loop ratio from
      uninterrupted runs}
\item One full-observation seed of the non-attentional prior was not
      reproducible across runs by a factor of $3.6$ at nominally
      identical configuration. \pend{resolve: the leading hypothesis is
      CPU thread count altering floating-point reduction order}
\end{itemize}

% =====================================================================
\section*{Code and data availability}
% =====================================================================
Implementation, verification harness, configurations and complete run
records are available at \pend{repository URL and archival DOI}. The
world model is a faithful small-scale DreamerV3; the quantum transformer
prior is a drop-in substitution installed on an already-constructed
agent, so the two arms differ in exactly one sublayer and the baseline
implementation is unmodified.

% =====================================================================
\section*{Acknowledgements}
% =====================================================================
\pend{acknowledgements, funding, hardware access}

% =====================================================================
\appendix
\section{Pending analysis}\label{app:remaining}
% =====================================================================
Every training run this paper requires is complete. The items below are
blocked on a single re-evaluation pass over the finished checkpoints.

\begin{enumerate}\itemsep3pt
\item Offline table, eight seeds, both arms, $k \in \{1,5,10,20\}$
      (Section~\ref{sec:offline}).
\item Context ablation, $K \in \{2,4,8,16\}$ at fixed width, eight seeds
      (Section~\ref{sec:context}).
\item \textbf{Done.} Closed-loop returns and the crossover table
      (Section~\ref{sec:closedloop}).
\item Optional strengthening: four further seeds per arm under partial
      observation, to move the main effect from marginally to
      comfortably powered.
\item Per-seed values and paired test for the open-loop comparison;
      full-observation row of Table~\ref{tab:imagination}
      (Section~\ref{sec:imagination}).
\item \textbf{Built.} Learning curves, circuit gradient trace and
      KL trajectories (\texttt{Paper\_B/make\_figures.py}). Verify the
      gradient column is non-empty before citing that figure.
\item Resolve the single irreproducible seed.
\item Novelty search record (Appendix~\ref{app:novelty}).
\item Repository URL, archival DOI, affiliation, acknowledgements, and
      full citation details for all references.
\end{enumerate}

% =====================================================================
\section{Novelty search}\label{app:novelty}
% =====================================================================
\begin{pendingbox}{search record}
The search has been carried out and cleared. Record here so the
first-of-its-kind claim is auditable: databases queried, exact search
strings, date of each query, number of records screened, and the
specific near misses with the reason each was excluded. The claim in the
abstract is only as strong as this appendix.
\end{pendingbox}

% =====================================================================
\section{Methodological notes}\label{app:notes}
% =====================================================================
Four errors during this work produced plausible but wrong numbers rather
than failures. We record them because each is a general hazard, and
because a reader of the main text cannot otherwise rule them out.

\begin{enumerate}\itemsep3pt
\item \textbf{A run-naming collision destroyed two completed runs.} An
      earlier scheme appended only non-default settings to the directory
      name, so an ablation at a default value produced the same name as
      a headline run at a different step budget and overwrote it. Both
      runs completed normally and left plausible result files. Fix: every
      setting appears unconditionally, and training refuses to write
      into a directory containing a completed run.
\item \textbf{A transformation mismatch inverted a comparison.} The
      open-loop measurement of Section~\ref{sec:imagination} initially
      compared decoder output against raw observations, but the decoder
      is trained against $\mathrm{symlog}(o)$. The mismatch applied an
      unmodelled transformation to one side and reversed the sign of the
      result. All reported figures are in symlog space.
\item \textbf{A confounded ablation.} The first context sweep varied
      window length and model width together, admitting no trend. It was
      discarded and re-run at fixed width. The naming guard prevents
      overwriting but not confounding, so the sweep driver now refuses
      to start when a result-affecting setting absent from the directory
      name differs from its default.
\item \textbf{An analysis grouping merged distinct conditions.} A paired
      comparison keyed on seed and arm alone, so every condition
      collided and whichever directory sorted last silently won.
      Grouping is now keyed on the full configuration, with architecture
      inferred from parameter count rather than from a configuration
      field that older runs lack.
\item \textbf{An analysis key omitted the environment.} The Phase~1
      aggregation grouped runs by model, capacity and horizon, but not
      by task. CartPole and Pendulum runs therefore collided on shared
      seeds, and because ``pendulum'' sorts after ``cartpole'' the
      Pendulum value silently replaced its CartPole counterpart. The
      effect was not subtle: Pendulum's classical one-step error is
      about $0.00425$ against CartPole's $0.00107$, so a single pooled
      run inflated the classical mean fourfold and produced an apparent
      $29\%$ quantum improvement ($0.711\times$, $p = 0.041$) where the
      true CartPole figure is $2.5\%$ ($0.975\times$, same $p$, same
      sign pattern). The $p$-value was unchanged throughout, because
      only the magnitude was spurious --- which is precisely why the
      error would have survived a significance check. Context length
      was missing from the same key, so three of every four runs at each
      width were discarded unseen.

\end{enumerate}

All five share a structure: a quantity that varied was not
recorded as having varied. Three of the five changed a reported
number and two were found only because a check was added that had not
previously existed.

We draw two lessons. First, run provenance deserves the same defensive
treatment as numerical code: a directory name is a cheap place to
record what varied, and an aggregation key is a cheap place to verify
it. Second, and less comfortably, a grouping key written by listing the
fields one is thinking about will omit the fields one is not. The
correct default is to key on the entire configuration and explicitly
exclude only the axes under comparison. Every aggregation reported in
this paper now does so, and halts with an error rather than selecting a
winner when two runs collide on a key --- the check that found item~5,
and that would have found items~1 and~3 when they occurred.

% =====================================================================
\bibliographystyle{unsrtnat}
\begin{thebibliography}{9}
\bibitem{hafner2023} D.~Hafner, J.~Pasukonis, J.~Ba, T.~Lillicrap,
  \emph{Mastering Diverse Domains through World Models},
  arXiv:2301.04104. \pend{venue of record}
\bibitem{transdreamer} C.~Chen, Y.-F.~Wu, J.~Yoon, S.~Ahn,
  \emph{TransDreamer: Reinforcement Learning with Transformer World
  Models}, arXiv:2202.09481 (2022). \pend{verify author list}
\bibitem{chen2020} S.~Y.-C.~Chen et al., \emph{Variational Quantum
  Circuits for Deep Reinforcement Learning}. \pend{full citation}
\bibitem{skolik2023} \emph{Variational Quantum Circuit Design for
  Quantum Reinforcement Learning on Continuous Environments},
  arXiv:2312.13798. \pend{verify authors}
\bibitem{dreamqas} \emph{DreamQAS: Learning a Decision-Useful World
  Model for VQE-Efficient Quantum Architecture Search},
  arXiv:2607.29491. \pend{verify authors}
\bibitem{qasa} \emph{Quantum Adaptive Self-Attention for Quantum
  Transformer Models}, arXiv:2504.05336. \pend{verify authors; cite the
  revised version}
\bibitem{schuld2021} M.~Schuld, R.~Sweke, J.~J.~Meyer, \emph{Effect of
  data encoding on the expressive power of variational quantum machine
  learning models}, Phys.\ Rev.\ A \textbf{103}, 032430 (2021).
\bibitem{cerezo2021} M.~Cerezo, A.~Sone, T.~Volkoff, L.~Cincio,
  P.~J.~Coles, \emph{Cost function dependent barren plateaus in shallow
  parametrized quantum circuits}, Nat.\ Commun.\ \textbf{12}, 1791
  (2021). \pend{verify}
\bibitem{bowles2024} J.~Bowles, S.~Ahmed, M.~Schuld, \emph{Better than
  classical? The subtle art of benchmarking quantum machine learning
  models}, arXiv:2403.07059 (2024).
\bibitem{papera} V.~Liutarevich, \emph{How Many Seeds? Dispersion,
  Resolving Power, and a Failed Replication in Quantum Sequence
  Modelling}. \pend{companion paper}
\end{thebibliography}

\end{document}
